% Part 1 closing chapter: the virial partition at every scale (first half of the Virial material; the cosmological coupling opens Part 2).
% Also: docs/papers/latex/virial_partition_atoms_to_horizon (numbers: docs/verification/scripts/verify_virial_atoms_to_horizon.py).
% Source: the five Virial papers (read in full 2026-10-02): Virial Efficiency (25 Feb 2026); Wide Domains (25 Feb 2026); Virial Partners (Mar 2026);
% Gravitational Decoherence and the Virial Partition (Mar 2026); The Thermodynamic Identity Governing the Virial Theorem (PRL version, 18 Mar 2026).
% Checks: docs/verification/virial/VIRIAL_CHECK.md, NBODY_TRACE.md.
\providecommand{\authorhold}[1]{\par\smallskip\noindent\fbox{\parbox{0.95\linewidth}{\small\textbf{Author hold.} #1}}\par\smallskip}

\chapter{One half at every scale: the virial partition}\label{ch:virial_law}

\section*{Overview}
IAM's law says that every irreversible transition from a quantum superposition to a classical record pays $k_BT\ln2$ per bit. This chapter shows where that
cost sits in the most widely tested equilibrium result in physics, the virial theorem of Clausius (1870). The theorem fixes the share of binding energy that a
stable system bound by a $1/r$ potential holds as motion: exactly one half. The same half holds from the hydrogen atom to clusters of galaxies, and a short
thermodynamic argument identifies it: the kinetic half is the irreversible cost of the binding transition itself. Part~2 opens by applying this to the universe,
where it fixes the one coupling that every cosmological result in this book carries.

\section{The theorem and why the factor is one half}
For a bound system whose potential energy is a homogeneous function of degree $k$ in the coordinates, the time averages of kinetic and potential energy satisfy
\begin{equation}
2\langle K\rangle=k\,\langle V\rangle .
\end{equation}
This follows from Euler's theorem for homogeneous functions applied to the virial $\sum_i\mathbf r_i\cdot\mathbf p_i$, whose time average is constant for a
bound system. Gravity and the Coulomb interaction both have $V\propto r^{-1}$, so $k=-1$ and
\begin{equation}
2\langle K\rangle+\langle V\rangle=0,\qquad \langle K\rangle=\tfrac12|\langle V\rangle| .
\end{equation}
The factor is set by the degree of the potential, and the degree is set by geometry: in three spatial dimensions Gauss's law makes the field of a point source
fall as $1/r^2$, so the potential falls as $1/r$. Nothing in the result depends on the number of bodies, the mass distribution or the scale. The same
$1/2$ appears in the Bohr atom, in a star, and in a cluster of galaxies.

\section{The same one half across physical scales}
Table~\ref{tab:virial_domains} collects the domains in which the $1/r$ virial ratio has been tested.

\begin{table}[h]\centering\small
\caption{The virial ratio for $1/r$ binding, from atoms to clusters of galaxies.}\label{tab:virial_domains}
\begin{tabular}{llll}
\hline
System & Scale & Quantity & Result\\\hline
Hydrogen atom & $10^{-10}$ m & $\langle K\rangle/|\langle V\rangle|$ & $1/2$ exactly (13.6\,eV of 27.2\,eV)\\
20 atoms, H--Xe (Hartree--Fock limit) & $10^{-10}$ m & $T/|V|$ & $1/2$ (exact for converged wavefunctions)\\
10 molecules, H$_2$--CO$_2$ & $10^{-9}$ m & $T/|V|$ & $1/2$ at equilibrium geometry\\
Thermal equilibrium & all & energy per quadratic d.o.f. & $k_BT/2$\\
White dwarfs & $10^{7}$ m & Chandrasekhar mass & $1.456\,M_\odot$ ($\mu_e=2$); none observed above it\\
The Sun & $10^{9}$ m & $|U|/2L$ & 24 Myr: contraction cannot power it, so it burns hydrogen\\
Galaxy clusters & $10^{23}$ m & dynamical vs lensing mass & agree to tens of percent (Zwicky 1933 onward)\\
Simulated halos, $10^{10}$--$10^{15}M_\odot$ & $10^{21}$--$10^{23}$ m & $2K/|W|$ & $1.1$--$1.3$; $1.02$--$1.17$ with surface pressure\\
Schwarzschild black hole & $r_s$ & $T_HS/Mc^2$ & $1/2$ (Smarr relation; Chapter~\ref{ch:blackholes})\\
Cosmic horizon & $10^{26}$ m & $\beta_m/\Omega_m$ & $1/2$ by prediction; consistent with Planck (Chapter~\ref{ch:virial})\\\hline
\end{tabular}
\end{table}

For converged Hartree--Fock wavefunctions the ratio is exact by construction: any wavefunction optimised under a uniform scaling of the coordinates satisfies
the quantum virial theorem. The atomic rows therefore show that quantum mechanics obeys the theorem; they are a statement of the theorem's reach, not an
independent measurement. The gravitational rows are measurements, at the precision the systems allow. The span from the atom to the cluster is 33 orders of
magnitude in size; Part~2 carries the same half to the horizon of a black hole and to the horizon of the universe.

\section{What the kinetic half is: the thermodynamic identity}
Clausius derived the theorem from Newton's second law. The derivation describes the bound state; it says nothing about the transition into it. Consider a
system that falls from rest at infinity ($E_i=0$) into a stable bound state with $E_f=\langle K\rangle+\langle V\rangle<0$, releasing the difference to its
surroundings.
\begin{enumerate}
\item First law: the heat released is $Q=E_i-E_f=-(\langle K\rangle+\langle V\rangle)>0$.
\item Second law: the transition is irreversible and produces entropy $\Delta S=Q/T$.
\item Landauer: the minimum cost of producing $\Delta S$ irreversibly is $E_{\rm L}=T\Delta S=Q$. The temperature cancels.
\item The $1/r$ equilibrium condition $2\langle K\rangle+\langle V\rangle=0$ gives $-E_f=\langle K\rangle$.
\end{enumerate}
Together,
\begin{equation}
\langle K\rangle=Q=T\Delta S=E_{\rm L}=\tfrac12|\langle V\rangle| .
\end{equation}
The hydrogen atom shows it directly: an electron captured from rest into the ground state emits a 13.6\,eV photon, the binding energy, which equals the
electron's mean kinetic energy in that state; the potential energy is $-27.2$\,eV. The energy carried away when a $1/r$ system binds equals its kinetic half.

This does not derive the virial theorem from thermodynamics; the fourth step is the mechanical condition. What it does is identify the kinetic half with a
definite thermodynamic quantity: the irreversible cost of the binding transition, carried off to the surroundings and recorded there as entropy at
$k_BT\ln2$ per bit. Jacobson (1995) obtained the Einstein equations from $\delta Q=T\,\delta S$ on local Rindler horizons; the same relation applied at the
boundary of a forming bound state returns the virial partition. The identity holds wherever the binding energy is released --- radiated by an atom, by a
contracting gas cloud, by a forming star. For collisionless dark matter, which relaxes without radiating, the energy is redistributed rather than released,
and the identity applies to the information written in the relaxation, not to emitted heat; this is the step IAM takes at cosmological scale.
